\documentclass[11pt]{article}
\usepackage[margin=1in]{geometry}
\usepackage{amsmath,amssymb,amsthm}
\usepackage[hidelinks]{hyperref}
\newtheorem{theorem}{Theorem}
\newtheorem{lemma}[theorem]{Lemma}
\newcommand{\qi}[1]{[#1]_q}
\title{Bottom heaviness of the four \(q\)-Stirling families\\
\large A solution to MathDB problem 354788}
\author{Principia Math in Progress}
\date{September 12, 2026}
\begin{document}
\maketitle

\paragraph{Task 1. Solve the original conjecture.}
For a nonzero polynomial \(F(q)=\sum_{i=0}^{D}a_iq^i\) of degree \(D\),
bottom heaviness means \(a_i\ge a_{D-i}\) for \(0\le i<D/2\).
We prove this for all four polynomials in Conjecture 7.6 of
Sagan--Swanson \cite{SS}. Constants and zero polynomials satisfy the
condition vacuously.

Put \(\qi{m}=1+q+\cdots+q^{m-1}\) for integers \(m\ge1\), and
\(\qi{0}=0\). The source's four families have initial values
\(S[0,k]=c[0,k]=S_B[0,k]=c_B[0,k]=\delta_{0k}\) and recurrences
\begin{align}
 S[n,k]&=S[n-1,k-1]+\qi{k}S[n-1,k], \label{eq:A2}\\
 c[n,k]&=c[n-1,k-1]+\qi{n-1}c[n-1,k], \label{eq:A1}\\
 S_B[n,k]&=S_B[n-1,k-1]+\qi{2k+1}S_B[n-1,k], \label{eq:B2}\\
 c_B[n,k]&=c_B[n-1,k-1]+\qi{2n-1}c_B[n-1,k]. \label{eq:B1}
\end{align}
Here \(n\ge1\), \(0\le k\le n\), and entries outside the triangular
range are zero. These are equations (2), (7), (12), and (14) of
\cite{SS}. In particular, \(S[n,k]\) is the unshifted version in (2).

\begin{lemma}\label{lem:product}
Every finite product of \(q\)-integers \(\qi{m}\), \(m\ge1\), has
nonnegative, symmetric, unimodal coefficients.
\end{lemma}
\begin{proof}
We give an elementary induction that also handles plateaux. Suppose
\(f(q)=\sum_{j=0}^{d}f_jq^j\) has the asserted properties, and set
\(f_j=0\) for all integers outside \([0,d]\).
The value \(f_j\) is a nonincreasing function of
\(\lvert j-d/2\rvert\) on these integer indices. For
\(g=f\qi{m}\), symmetry follows by reversing both factors. Moreover,
\[
 g_i-g_{i-1}=f_i-f_{i-m}.
\]
If \(1\le i\le(d+m-1)/2\), then
\[
 (i-m-d/2)^2-(i-d/2)^2=m(d+m-2i)\ge0.
\]
Thus \(f_i\ge f_{i-m}\), so \(g\) is nondecreasing up to its center
\((d+m-1)/2\). Symmetry gives unimodality. Induction starts with the
empty product \(1\).
\end{proof}

\begin{lemma}\label{lem:sum}
A finite nonnegative linear combination of nonzero polynomials
having nonnegative, symmetric, unimodal coefficients is bottom heavy.
\end{lemma}
\begin{proof}
Discard zero weights. If all weights vanish the assertion is vacuous;
otherwise let \(D\) be the largest degree of a summand. Nonnegativity
ensures that \(D\) is the degree of the sum.
Fix a summand \(f\) of degree \(d\le D\), extend its coefficient
sequence by zeros, and let \(0\le i<D/2\). Since
\[
 (D-i-d/2)^2-(i-d/2)^2=(D-2i)(D-d)\ge0,
\]
the same distance-from-center observation gives \(f_i\ge f_{D-i}\).
Multiplying by the nonnegative weight and adding over the summands
proves the claim.
\end{proof}

\begin{theorem}\label{thm:main}
For every admissible \(n,k\), the coefficient sequences of
\(S[n,k]\), \(c[n,k]\), \(S_B[n,k]\), and \(c_B[n,k]\) are bottom heavy.
More generally, every finite nonnegative linear combination of
products of \(q\)-integers is bottom heavy.
\end{theorem}
\begin{proof}
The general assertion is the combination of Lemmas
\ref{lem:product} and \ref{lem:sum}. Let \(\mathcal C\) consist of
these linear combinations, including \(0\) and the empty product \(1\).
This class is closed under addition and multiplication by any
\(\qi{m}\), \(m\ge0\). Every coefficient multiplying an entry in
\eqref{eq:A2}--\eqref{eq:B1} is such a \(q\)-integer.
Induction on \(n\) therefore places every entry in \(\mathcal C\).
This also handles \(k=0\), \(k=n\), and \(n=0\) without a positivity
or degree assumption on a zero entry.
\end{proof}

The comparison in Lemma \ref{lem:sum} uses each summand's own center:
the summands need not have equal degrees. For example
\(S[4,3]=\qi1+\qi2+\qi3=3+2q+q^2\) has the required asymmetry.
The theorem settles the entire original conjecture; it does not
assert unimodality of a sum with different centers.

\paragraph{Verification.}
The proof is independent of computation. The accompanying exact
integer checker reconstructs all four recurrences for \(n\le50\),
compares them with explicit complete-homogeneous and elementary
symmetric-function expansions for \(n\le10\), and tests the lemmas
on products and mixtures with unequal degrees. It also reproduces
the source's nonunimodal type \(B\) examples.

\paragraph{Task 2. Construct coefficient injections.}
Translate the summand comparison into an explicit injection from
objects counted by the coefficient of \(q^{D-i}\) to those counted
by \(q^i\), using products of chains and the source's partition and
permutation models.

\paragraph{Task 3. Classify equality and quantify the asymmetry.}
Determine exactly when \(a_i=a_{D-i}\) in each family, and obtain
effective lower bounds for \(a_i-a_{D-i}\) outside those cases.

\paragraph{Task 4. Extend the theorem to further Stirling analogues.}
Identify which \(q\)-Stirling polynomials for complex reflection
groups have nonnegative expansions in products of \(q\)-integers,
then apply the theorem and characterize any exceptions.

\paragraph{Task 5. Resolve the stronger coefficient conjectures.}
Prove or refute coefficient log-concavity in type \(A\) and parity
log-concavity in type \(B\), as stated in Conjectures 7.4 and 7.5 of
\cite{SS}; isolate additional structure beyond the sum lemma.

\begin{thebibliography}{9}
\bibitem{SS}
B.~E. Sagan and J.~P. Swanson,
\emph{\(q\)-Stirling numbers in type \(B\)},
European Journal of Combinatorics \textbf{118} (2024), 103899.
\href{https://par.nsf.gov/servlets/purl/10626831}{Published text};
\href{https://arxiv.org/abs/2205.14078}{arXiv:2205.14078}.
Definitions: pp.~2--4; Conjecture 7.6: p.~31.
\end{thebibliography}
\end{document}
